\documentclass[11pt]{article}
\usepackage{amsmath, amssymb, amsthm, amsfonts}
\usepackage{mathtools}
\usepackage{bm}
\usepackage{geometry}
\geometry{margin=1in}

\title{Prime–Indexed Quantum Zeta Dynamics and Multiplicity–Stabilized Two–Qubit Models:\\
Technical Report and Defensive Prior–Art Publication}
\author{Citizen Gardens -- The Foundation of Multiplicity\\
The Prime Materia Commons\\
\texttt{info@citizengardens.org}}
\date{May 14, 2026}

\begin{document}
\maketitle

\begin{abstract}

This document records a sequence of constructions and corrections relating to prime–indexed dynamical systems motivated by Multiplicity Theory. The main components are:

\begin{itemize}
  \item A zeta–driven quantum dynamical system on a graded Hilbert space, the Zeta–Multiplicity Oscillator (ZMO), with evolution generated by an operator of the form
  \(\mathcal{U}_{\rm ZMO} = \mathbf{A}_{\rm prime} + \mathbf{B}_{\rm zeta} + \mathbf{E}_{\rm internal} + \Lambda_m \mathbf{\Xi}_{\rm bridge}\),
  where the prime sector and the zeta sector are coupled using explicit–formula style terms.
  \item A series of minimal executable kernels for a ``Universal Multiplicity Constant'' \(\Lambda_m\) acting as a stabilizer in two–qubit models of neurotransmitter–like subsystems, culminating in a matched–control kernel in which the stabilization claim is falsified.
  \item A clarified separation between (i) formal operator–level proposals (ZMO, PIRTM), (ii) concrete finite–dimensional toy models, and (iii) empirically falsifiable numerical effects.
\end{itemize}

The defensive content is as follows:

\begin{enumerate}
  \item Definition of the Zeta–Multiplicity Oscillator on a prime–graded Hilbert space, with an explicit operator decomposition and constraints required for boundedness, self–adjointness, and control of the essential spectrum. The construction references known connections between Riemann zeros and quantum chaos/random matrix spectra.\cite{RZQC,RMQChaos,RHOverview}
  \item Formalization of \(\Lambda_m\) as a bounded, prime–indexed scalar functional \(\Lambda_m^{\mathrm{op}}(t)\) constructed from a golden–ratio skeleton \(\Phi^{p_i}\) and prime–frequency oscillations, and insertion of this functional as a coupling factor in two–qubit Hamiltonians.
  \item A complete, runnable two–qubit kernel with:
  \begin{enumerate}
    \item Hilbert space \(\mathbb{C}^2 \otimes \mathbb{C}^2\),
    \item maximally entangled initial state \((|00\rangle + |11\rangle)/\sqrt{2}\),
    \item ZZ coupling scaled either by \(\Lambda_m^{\mathrm{op}}(t)\) or by its time average,
    \item prime–parameterized local rotations,
    \item concurrence as entanglement observable,
    \item and a matched scalar control with identical mean coupling strength.
  \end{enumerate}
  \item An explicit negative result: for this Hamiltonian and observable, the prime–oscillatory coupling produces \emph{more} entanglement decay than the constant coupling of matched average strength. The initial stabilization claim is therefore retracted for this toy model.
\end{enumerate}

All constructions are recorded in sufficient technical detail to constitute prior art for operator forms, scalar functionals, and finite–dimensional kernels coupling primes, golden–ratio structures, and quantum dynamics.
\end{abstract}
\section{Zeta–Multiplicity Oscillator (ZMO)}

\subsection{Ambient Hilbert Space and State}

Let \(\mathbb{P}\) denote the set of prime numbers. Consider the graded Hilbert space
\[
  \mathcal{H}
  = \ell^2(\mathbb{P}) \,\widehat{\otimes}\, L^2(\mathbb{R}) \,\widehat{\otimes}\, \mathbb{C}^d,
\]
where:
\begin{itemize}
  \item \(\ell^2(\mathbb{P})\) carries a canonical orthonormal basis \(\{ e_p : p \in \mathbb{P} \}\), representing prime–indexed modes.
  \item \(L^2(\mathbb{R})\) carries a basis adapted to zeta zeros (e.g.\ wavepackets built from eigenfunctions of a multiplication operator in the spectral representation).
  \item \(\mathbb{C}^d\) accounts for internal degrees of freedom (spin, internal registers, or auxiliary multiplicity channels).
\end{itemize}

A generic oscillator state is written as
\begin{equation}
  |\Psi_{\rm osc}(t)\rangle
  = \sum_{p\in\mathbb{P}} \alpha_p \, e^{i \omega_p(t) t} \,|\psi_p\rangle
    \,\otimes\, \sum_{\rho} \alpha_{\rho} \, e^{i \rho t}|\psi_{\rho}\rangle,
\end{equation}
where:
\begin{itemize}
  \item \(\rho\) ranges over nontrivial zeros of the Riemann zeta function \(\zeta(s)\).\cite{RHOverview,RZQC}
  \item \(\omega_p(t)\) is a prime–encoded frequency, defined below.
  \item \(|\psi_p\rangle\) and \(|\psi_{\rho}\rangle\) are normalized vectors in the respective subspaces.
\end{itemize}

\subsection{Prime–Encoded Frequency via Euler Product}

For \(\Re(s) > 1\), the Euler product representation
\begin{equation}
  \zeta(s) = \prod_{p} \frac{1}{1 - p^{-s}}
\end{equation}
implies that the logarithmic derivative \(\zeta'(s)/\zeta(s)\) can be expressed as a convergent prime sum.\cite{RHOverview}
Motivated by this structure, a prime–encoded effective frequency is defined by
\begin{equation}
  \omega_p(t)
  = \sum_{p'} \frac{\alpha_{p'}}{{p'}^s} \, t,
\end{equation}
for some fixed \(s\) with \(\Re(s) > 1\), and fixed complex coefficients \(\alpha_{p'}\) satisfying summability conditions to ensure convergence.

This construction realizes a frequency sector in which each prime mode carries a time–dependent phase factor that is a linear functional of the Euler–product weights.

\subsection{Zeta–Multiplicity Operator}

Define the Zeta–Multiplicity Operator
\begin{equation}
  \mathcal{U}_{\rm ZMO}
  = \mathbf{A}_{\rm prime} + \mathbf{B}_{\rm zeta} + \mathbf{E}_{\rm internal} + \Lambda_m \,\mathbf{\Xi}_{\rm bridge}.
\end{equation}

\subsubsection{Prime Sector \(\mathbf{A}_{\rm prime}\)}

The prime sector is given by
\begin{equation}
  \mathbf{A}_{\rm prime} = D_\sigma + K,
\end{equation}
where:
\begin{itemize}
  \item \(D_\sigma\) is a diagonal (or multiplication) operator acting in \(\ell^2(\mathbb{P})\), e.g.
  \(D_\sigma e_p = \sigma(p) e_p\), where \(\sigma:\mathbb{P}\to\mathbb{R}\) is chosen so that
  \(D_\sigma\) is bounded (or essentially self–adjoint on a dense domain).
  \item \(K\) is a compact (or Hilbert–Schmidt) operator, with \(\alpha > 1/2\) in a Schatten–class constraint ensuring \(K\) is Hilbert–Schmidt.\cite{RZQC}
\end{itemize}

\subsubsection{Zeta Sector \(\mathbf{B}_{\rm zeta}\)}

The zeta sector \(\mathbf{B}_{\rm zeta}\) acts on \(L^2(\mathbb{R})\) as a Fourier multiplier operator:
\begin{equation}
  (\widehat{\mathbf{B}_{\rm zeta} f})(\xi)
  = m(\xi) \,\hat{f}(\xi),
\end{equation}
where the symbol \(m(\xi)\) is constructed from zeta–related quantities, for example:
\begin{itemize}
  \item \(m(\xi) = \Re \log \zeta(\tfrac12 + i\xi)\), or
  \item a smoothed version involving a Fejér kernel localized around windows of height \(T\),
\end{itemize}
subject to the requirement that \(m \in L^\infty(\mathbb{R})\) to ensure boundedness and, in the self–adjoint case, that \(m(\xi)\) is real–valued almost everywhere.\cite{RZQC,KeatingQC}

\subsubsection{Internal Sector \(\mathbf{E}_{\rm internal}\)}

The internal sector \(\mathbf{E}_{\rm internal}\) is a self–adjoint operator on \(\mathbb{C}^d\), typically finite–dimensional, used to encode auxiliary multiplicity degrees of freedom.

\subsubsection{Bridge Sector \(\mathbf{\Xi}_{\rm bridge}\)}

The bridge \(\mathbf{\Xi}_{\rm bridge}\) couples prime and zeta sectors via terms reminiscent of the explicit formula:
\begin{equation}
  (\mathbf{\Xi}_{\rm bridge} \Psi)(p,\rho,\cdot)
  = \sum_{p_i} \sum_{\rho_k} \kappa(p_i,\rho_k) \cos(\gamma_k \log p_i) \,\Pi_{p_i, \rho_k}\Psi,
\end{equation}
where:
\begin{itemize}
  \item \(\rho_k = \tfrac12 + i\gamma_k\) range over nontrivial zeros of \(\zeta\).\cite{RHOverview}
  \item \(\kappa\) is a coupling kernel, and
  \item \(\Pi_{p_i,\rho_k}\) are projections onto modes labeled by \((p_i,\rho_k)\).
\end{itemize}
To avoid uncontrolled contributions from zero crowding, \(\kappa\) is taken to be localized using band–limited test functions (e.g.\ Fejér kernels), so that only finitely many \(\rho_k\) contribute significantly in each spectral window.\cite{RZQC}

\subsubsection{Multiplicity Constant \(\Lambda_m\) and Contractivity}

A scalar multiplicity constant \(\Lambda_m \in (0,1)\) is introduced to enforce contractivity constraints:
\begin{equation}
  \|\Phi_t\| \leq 1 - \varepsilon,
\end{equation}
where \(\Phi_t\) is the evolution family associated with \(\mathcal{U}_{\rm ZMO}\) in the semigroup or unitary picture.

In the unitary case, a Cayley transform
\begin{equation}
  U = (\mathcal{U}_{\rm ZMO} - iI)(\mathcal{U}_{\rm ZMO} + iI)^{-1}
\end{equation}
is used to map a self–adjoint operator to a unitary operator, while in the contractive case one directly considers the generator \(\mathcal{A} = -\mathcal{U}_{\rm ZMO}\) of a contraction semigroup \(e^{t\mathcal{A}}\).

\subsection{Essential Spectrum and Random–Matrix Analogies}

Under suitable conditions (bounded \(\mathbf{B}_{\rm zeta}\), compact \(\mathbf{\Xi}_{\rm bridge}\)), the essential spectrum satisfies
\begin{equation}
  \sigma_{\rm ess}(\mathcal{U}_{\rm ZMO})
  = \sigma_{\rm ess}(\mathbf{A}_{\rm prime}) + \operatorname{EssRan}(m) + \sigma_{\rm ess}(\mathbf{E}_{\rm internal}),
\end{equation}
where \(\operatorname{EssRan}(m)\) is the essential range of the symbol \(m\).\cite{RZQC}

Connections to quantum chaos and random matrix theory arise by comparing the local statistics of the spectrum of (a discretized version of) \(\mathcal{U}_{\rm ZMO}\) to those of classical random matrix ensembles (e.g.\ GUE), in line with known correspondence between Riemann zeros and eigenvalues of chaotic Hamiltonians.\cite{RZQC,KeatingQC,RandomMatrixQChaos}

\section{Prime–Indexed Multiplicity Functional \texorpdfstring{$\Lambda_m^{\mathrm{op}}(t)$}{Lambda\_m^{op}(t)}}

\subsection{Definition}

Fix a finite set of primes \(\{p_i\}_{i=1}^n\), a golden–ratio constant
\[
  \Phi = \frac{1 + \sqrt{5}}{2},
\]
and define a skeleton
\[
  \xi(p_i) = \Phi^{p_i}.
\]

For a real time parameter \(t \ge 0\), define the multiplicity functional
\begin{equation}
  \Lambda_m^{\mathrm{op}}(t)
  = \frac{1}{2} \,\frac{\sum_{i=1}^n \xi(p_i)\cos(p_i t)}{1 + \bigl|\sum_{i=1}^n \xi(p_i)\cos(p_i t)\bigr|}.
\end{equation}
By construction, \(|\Lambda_m^{\mathrm{op}}(t)| \le \tfrac12\) for all \(t\), and \(\Lambda_m^{\mathrm{op}}(t)\) is a bounded, prime–indexed scalar functional encoding prime frequencies and golden–ratio scaling.

\subsection{Time Average}

For a fixed time horizon \(T > 0\) and integration step \(\Delta t\), the time average is approximated by
\begin{equation}
  \overline{\Lambda_m^{\mathrm{op}}}
  = \frac{1}{N}\sum_{k=0}^{N-1} \Lambda_m^{\mathrm{op}}(t_k),
  \quad t_k = k \Delta t, \quad N \approx T/\Delta t.
\end{equation}
This average is used to construct a scalar control coupling in the matched–control kernel.

\section{Two–Qubit Neurotransmitter Kernel}

\subsection{Hilbert Space and Initial State}

Consider the two–qubit Hilbert space
\[
  \mathcal{H}_{\rm NT} = \mathbb{C}^2 \otimes \mathbb{C}^2,
\]
with computational basis \(\{|00\rangle,|01\rangle,|10\rangle,|11\rangle\}\).

Interpret the first qubit as a ``dopamine'' subsystem and the second as a ``serotonin'' subsystem. The initial state is the maximally entangled Bell–type state
\begin{equation}
  |\psi_0\rangle = \frac{1}{\sqrt{2}}(|00\rangle + |11\rangle).
\end{equation}

\subsection{Prime–Parameterized Local Drives}

Let \(p_1, p_2\) be distinct primes (for concreteness, \(p_1=2\), \(p_2=3\)). For time \(t\), define rotation angles
\begin{equation}
  \theta_{\rm dop}(t) = p_1 t, \qquad \theta_{\rm ser}(t) = p_2 t.
\end{equation}
Define local \(\sigma_y\) rotations on each qubit by
\begin{equation}
  R(\theta)
  = \begin{pmatrix}
      \cos\theta & -i\sin\theta \\
      -i\sin\theta & \cos\theta
    \end{pmatrix}.
\end{equation}
Embed these into the two–qubit space:
\begin{equation}
  U_{\rm dop}(t) = R(\theta_{\rm dop}(t)) \otimes I_2, \qquad
  U_{\rm ser}(t) = I_2 \otimes R(\theta_{\rm ser}(t)).
\end{equation}

Define the local Hamiltonian as the skew–Hermitian part of these unitaries:
\begin{equation}
  H_{\rm local}(t)
  = -\frac{i}{2}\bigl(U_{\rm dop}(t) - U_{\rm dop}(t)^\dagger\bigr)
    -\frac{i}{2}\bigl(U_{\rm ser}(t) - U_{\rm ser}(t)^\dagger\bigr).
\end{equation}
This ensures that \(H_{\rm local}(t)\) is Hermitian.

\subsection{ZZ Coupling Modulated by \texorpdfstring{$\Lambda_m^{\mathrm{op}}(t)$}{Lambda\_m^{op}(t)}}

Define the ZZ interaction operator
\begin{equation}
  Z = \begin{pmatrix}1 & 0\\ 0 & -1\end{pmatrix}, \qquad
  ZZ = Z \otimes Z
  = \operatorname{diag}(1,-1,-1,1).
\end{equation}

For each time \(t\), define the interaction Hamiltonian
\begin{equation}
  H_{\rm int}^{\rm prime}(t) = \Lambda_m^{\mathrm{op}}(t) \, ZZ,
\end{equation}
so that the total Hamiltonian in the prime–weighted run is
\begin{equation}
  H^{\rm prime}(t) = H_{\rm local}(t) + H_{\rm int}^{\rm prime}(t).
\end{equation}

For the scalar control, replace \(\Lambda_m^{\mathrm{op}}(t)\) by the time average \(\overline{\Lambda_m^{\mathrm{op}}}\):
\begin{equation}
  H_{\rm int}^{\rm scalar}(t) = \overline{\Lambda_m^{\mathrm{op}}} \, ZZ,
  \quad
  H^{\rm scalar}(t) = H_{\rm local}(t) + H_{\rm int}^{\rm scalar}(t).
\end{equation}
Thus, the only difference between the two Hamiltonians is whether the ZZ coupling is time–dependent and prime–structured, or constant with identical mean strength.

\subsection{Time Evolution}

For a small time step \(\Delta t\), approximate the unitary evolution over \([t,t+\Delta t]\) by
\begin{equation}
  U^{(\cdot)}(t,\Delta t)
  = \exp\bigl(-i H^{(\cdot)}(t) \Delta t\bigr),
\end{equation}
where \((\cdot)\) is either ``prime'' or ``scalar''. The state is updated via
\begin{equation}
  |\psi_{k+1}^{(\cdot)}\rangle
  = \frac{U^{(\cdot)}(t_k,\Delta t)|\psi_k^{(\cdot)}\rangle}{\|U^{(\cdot)}(t_k,\Delta t)|\psi_k^{(\cdot)}\rangle\|},
\end{equation}
with \(t_k = k \Delta t\), \(k=0,1,\dots,N-1\), \(N \approx T/\Delta t\).

\subsection{Observable: Concurrence}

For a pure two–qubit state
\[
  |\psi\rangle = a|00\rangle + b|01\rangle + c|10\rangle + d|11\rangle,
\]
the concurrence is defined by\footnote{Standard Wootters concurrence for pure two–qubit states.}
\begin{equation}
  C(\psi) = 2|ad - bc|.
\end{equation}
This is an entanglement–sensitive observable taking values in \([0,1]\).

\subsection{Minimal Executable Kernel: Exact Code}

The following Python kernel implements the above construction with three primes \(\{2,3,5\}\), computes the prime–weighted run, estimates the time average \(\overline{\Lambda_m^{\mathrm{op}}}\), and then executes a scalar control run using the matched average coupling. The code uses \texttt{numpy} and \texttt{scipy.linalg.expm} for matrix exponentials.

\medskip
\noindent\textbf{Code Snippet.}
\begin{verbatim}
import numpy as np
from scipy.linalg import expm

primes = np.array([2., 3., 5.])
phi = (1 + np.sqrt(5)) / 2
xi = phi ** primes

def Lambda_m_op(t, scalar_control=False, avg_lm=None):
    if scalar_control and avg_lm is not None:
        return avg_lm
    lm = 0.0
    for i, p in enumerate(primes):
        lm += xi[i] * np.cos(p * t)
    return lm / (1 + np.abs(lm)) * 0.5

def run_simulation(scalar_control=False, dt=0.05, T=5.0, avg_lm=None):
    times = np.arange(0, T + dt/2, dt)
    psi0 = np.array([1/np.sqrt(2), 0, 0, 1/np.sqrt(2)], dtype=complex)
    psi0 /= np.linalg.norm(psi0)
    states = [psi0.copy()]
    concurrences = []
    lm_values = []
    for idx in range(1, len(times)):
        t = times[idx-1]
        lm = Lambda_m_op(t, scalar_control, avg_lm)
        lm_values.append(lm)
        ZZ = np.kron(np.diag([1., -1.]), np.diag([1., -1.]))
        Hint = lm * ZZ
        theta_d = primes[0] * t
        theta_s = primes[1] * t
        R_d = np.array([[np.cos(theta_d), -1j*np.sin(theta_d)],
                        [-1j*np.sin(theta_d), np.cos(theta_d)]],
                       dtype=complex)
        R_s = np.array([[np.cos(theta_s), -1j*np.sin(theta_s)],
                        [-1j*np.sin(theta_s), np.cos(theta_s)]],
                       dtype=complex)
        U_d = np.kron(R_d, np.eye(2))
        U_s = np.kron(np.eye(2), R_s)
        H_local = -1j * (U_d - U_d.conj().T)/2 \
                  -1j * (U_s - U_s.conj().T)/2
        H = Hint + H_local
        U = expm(-1j * H * dt)
        psi_new = U @ states[-1]
        psi_new /= np.linalg.norm(psi_new)
        states.append(psi_new)
        a, b, c, d = psi_new
        conc = 2 * np.abs(a * d - b * c)
        concurrences.append(conc)
    return np.mean(concurrences), np.min(concurrences), np.mean(lm_values)

# Prime-weighted run (computes true average)
mean_prime, min_prime, avg_lm = run_simulation(False)

# Matched scalar control
mean_scalar, min_scalar, _ = run_simulation(True, avg_lm=avg_lm)

print("Prime-weighted run:")
print(f"  Mean concurrence: {mean_prime:.4f}")
print(f"  Min concurrence: {min_prime:.4f}")
print(f"  Avg Λ_m^op(t): {avg_lm:.4f}")

print("\nScalar control (matched average):")
print(f"  Mean concurrence: {mean_scalar:.4f}")
print(f"  Min concurrence: {min_scalar:.4f}")
print(f"  Constant Λ_m: {avg_lm:.4f}")
\end{verbatim}

For a representative execution with \(T=5.0\), \(\Delta t = 0.05\), the output is:
\begin{align*}
  \text{Prime-weighted run:} & \quad \text{Mean concurrence } \approx 0.9181, \quad
  \text{Min concurrence } \approx 0.7530, \quad
  \overline{\Lambda_m^{\mathrm{op}}(t)} \approx -0.0137;\\
  \text{Scalar control (matched):} & \quad \text{Mean concurrence } \approx 0.9985, \quad
  \text{Min concurrence } \approx 0.9951, \quad
  \Lambda_m = \overline{\Lambda_m^{\mathrm{op}}(t)} \approx -0.0137.
\end{align*}

\subsection{Result: Falsification of Prime–Indexed Stabilization for This Kernel}

For the above Hamiltonian and observable, with matched average coupling strength, the scalar control preserves concurrence near unity, while the prime–modulated coupling yields lower mean and minimum concurrence. Thus, in this toy model:
\begin{equation}
  C_{\rm scalar}(t) \ge C_{\rm prime}(t) \quad \text{for most } t \in [0,T],
\end{equation}
and in particular,
\begin{equation}
  \mathbb{E}_t[C_{\rm scalar}(t)] > \mathbb{E}_t[C_{\rm prime}(t)], \qquad
  \inf_t C_{\rm scalar}(t) > \inf_t C_{\rm prime}(t).
\end{equation}

Consequently, the initial claim that prime–indexed time dependence of \(\Lambda_m^{\mathrm{op}}(t)\) stabilizes concurrence relative to a scalar coupling of comparable strength is \emph{falsified} for this Hamiltonian and observable.

This negative result is itself part of the prior art: it demonstrates that naive insertion of a bounded prime–oscillatory \(\Lambda_m^{\mathrm{op}}(t)\) into a ZZ coupling does not automatically realize the desired contractivity and stabilization properties.

\section{Prior–Art Claims}

The developments above support the following prior–art statements:

\begin{enumerate}
  \item The definition of a Zeta–Multiplicity Oscillator on
  \(\mathcal{H} = \ell^2(\mathbb{P}) \widehat{\otimes} L^2(\mathbb{R}) \widehat{\otimes} \mathbb{C}^d\),
  with evolution generated by an operator of the form
  \(\mathcal{U}_{\rm ZMO} = \mathbf{A}_{\rm prime} + \mathbf{B}_{\rm zeta} + \mathbf{E}_{\rm internal} + \Lambda_m \mathbf{\Xi}_{\rm bridge}\),
  where \(\mathbf{B}_{\rm zeta}\) is a zeta–derived Fourier multiplier and \(\mathbf{\Xi}_{\rm bridge}\) is constructed using explicit–formula terms \(\cos(\gamma_k \log p_i)\), localized by band–limited kernels to control multi–zero contributions.\cite{RZQC,KeatingQC,RHOverview}
  \item The construction of a bounded, prime–indexed multiplicity functional
  \(\Lambda_m^{\mathrm{op}}(t)\) of the form
  \[
    \Lambda_m^{\mathrm{op}}(t)
    = \frac{1}{2}\frac{\sum_i \Phi^{p_i}\cos(p_i t)}{1 + \bigl|\sum_i \Phi^{p_i}\cos(p_i t)\bigr|},
  \]
  and its insertion into two–qubit Hamiltonians as a global stabilizing factor.
  \item The specific two–qubit neurotransmitter kernel described in Section~3, including:
  \begin{itemize}
    \item prime–parameterized local \(\sigma_y\) drives,
    \item ZZ coupling scaled either by \(\Lambda_m^{\mathrm{op}}(t)\) or by its time average,
    \item concurrence as entanglement observable,
    \item and a matched scalar control.
  \end{itemize}
  \item The explicit negative finding that, for this model, the prime–oscillatory stabilizer \(\Lambda_m^{\mathrm{op}}(t)\) induces stronger entanglement decay than the matched constant stabilizer.
\end{enumerate}

These items collectively define a technical region of prior art at the intersection of prime–indexed operator frameworks, zeta–driven quantum dynamics, golden–ratio–based functionals, and multiplicity–modulated two–qubit Hamiltonians.

\section{Prime-Indexed Quantum Teleportation and Arithmetic Coherence}

In this section, quantum teleportation is recast as a prime-indexed coherence test between arithmetic functors. The construction proceeds in three layers: (i) a cohomological formulation of teleportation, (ii) a prime-indexed stabilizer tableau representation, and (iii) an arithmetic refinement in which teleportation fidelity becomes a multiplicative functional on Pell- and Dirichlet-type data. A final subsection couples this coherence functional into the Genesis persistence dynamics as a Tele--Genesis interface.

\subsection{Cohomological Framing of Teleportation}

Consider a standard one-qubit teleportation protocol between Alice and Bob. Let $G$ be the underlying graph of the protocol, with vertices representing systems and edges representing entangled links. A deformation of the shared Bell resource is encoded by a $U(1)$-valued $1$-cochain
\[
\Phi \in C^{1}(G;U(1)),
\]
assigning to each entangled edge $e$ a phase $\Phi(e) \in U(1)$.

Bob's local correction operations (phase-undo gates) define a $U(1)$-valued $0$-cochain
\[
\Psi \in C^{0}(G;U(1)),
\]
supported on his vertex. The effective channel implemented by the protocol is the identity if and only if the edge phases induced by $\Phi$ are precisely cancelled by the coboundary of $\Psi$:
\[
\delta \Psi = \Phi,
\]
where $\delta: C^{0} \to C^{1}$ is the usual coboundary operator. In this language, ordinary teleportation with an undeformed Bell pair corresponds to the trivial cohomology class $[\Phi]=0$, and any residual deformation class $[\Phi] \neq 0$ manifests as a non-trivial unitary error on the teleported state.

\subsection{Prime-Indexed Stabilizer Tableau}

Now fix a finite set of prime labels $\mathcal{P}$ and consider a family of prime-indexed, deformed Bell states. For each prime $p \in \mathcal{P}$ let
\[
|\Phi\rangle^{(p)} = \frac{1}{\sqrt{2}}\bigl(|00\rangle + e^{i\phi_p}|11\rangle\bigr)
\]
denote a phase-deformed Bell pair shared between systems $B$ and $C$, with prime-indexed deformation phase $\phi_p \in \mathbb{R}/2\pi\mathbb{Z}$.

In stabilizer tableau form, a two-qubit generator row is written as
\[
(x_B,x_C \mid z_B,z_C \mid \varphi),
\]
where $(x_i,z_i)\in\mathbb{F}_2$ encode the Pauli string and $\varphi\in\mathbb{R}/2\pi\mathbb{Z}$ is the eigenvalue phase in $S|\psi\rangle = e^{i\varphi}|\psi\rangle$.

For the undeformed Bell pair $|\Phi^+\rangle = \frac{1}{\sqrt{2}}(|00\rangle+|11\rangle)$, a standard stabilizer generating set is
\[
\begin{array}{c|cccc|c}
\text{gen} & x_B & x_C & z_B & z_C & \varphi \\ \hline
g_1 & 1 & 1 & 0 & 0 & 0 \\
g_2 & 0 & 0 & 1 & 1 & 0
\end{array}.
\]
For the prime-deformed state $|\Phi\rangle^{(p)}$ one checks directly that
\[
Z_B Z_C |\Phi\rangle^{(p)} = + |\Phi\rangle^{(p)}, \qquad
X_B X_C |\Phi\rangle^{(p)} = e^{i\phi_p} |\Phi\rangle^{(p)},
\]
so its tableau is
\[
\begin{array}{c|cccc|c}
\text{gen} & x_B & x_C & z_B & z_C & \varphi \\ \hline
g_1^{(p)} & 1 & 1 & 0 & 0 & \phi_p \\
g_2        & 0 & 0 & 1 & 1 & 0
\end{array}.
\]
The deformation $1$-cochain $\Phi(p)=e^{i\phi_p}$ appears precisely as the phase entry on the $X_B X_C$ row; the prime $p$ itself is an external label on this tableau block and is never altered by Clifford propagation.

\subsection{\texorpdfstring{$n$}{n}-Qubit Prime-Indexed Tableau Teleportation}

Let $\vec p = (p_1,\dots,p_n)$ be an $n$-tuple of primes and let $\vec\phi = (\phi_{p_1},\dots,\phi_{p_n})$ be the corresponding deformation phases. Consider $2n$ qubits grouped into pairs $(B_k,C_k)$, and define the prime-indexed resource
\[
|\Phi\rangle^{(\vec p)} = \bigotimes_{k=1}^n |\Phi\rangle^{(p_k)}_{B_k C_k}
= \bigotimes_{k=1}^n \frac{1}{\sqrt{2}}\bigl(|00\rangle + e^{i\phi_{p_k}}|11\rangle\bigr)_{B_k C_k}.
\]
For each block $k$, generators are
\[
g_1^{(p_k)} = X_{B_k} X_{C_k}, \quad \varphi=\phi_{p_k},\qquad
g_2^{(p_k)} = Z_{B_k} Z_{C_k}, \quad \varphi=0.
\]
The full $2n$-qubit stabilizer tableau is block-diagonal in the $(x\mid z)$ part, with phase column
\[
\vec\varphi_{\mathrm{res}} =
(\underbrace{\phi_{p_1},0}_{\text{block }1},
 \underbrace{\phi_{p_2},0}_{\text{block }2},
 \dots,
 \underbrace{\phi_{p_n},0}_{\text{block }n})^{\mathsf T}.
\]
Each prime $p_k$ thus indexes an independent tableau block; the phase entry $\phi_{p_k}$ in the $X_{B_k} X_{C_k}$ row is the only place where the deformation appears.

\paragraph{Protocol.} Alice holds input qubits $A_1,\dots,A_n$. For each $k$ she performs the standard teleportation sequence:
\begin{itemize}
  \item CNOT from $A_k$ to $B_k$, then Hadamard on $A_k$,
  \item measurement of $(A_k,B_k)$ in the computational basis, yielding bits $(a_k,b_k)\in\{0,1\}^2$.
\end{itemize}
Bob holds $C_1,\dots,C_n$ and applies, on each qubit, either:
\begin{align*}
\text{Pauli-only correction:} \quad & V_{a_k,b_k} = X^{b_k} Z^{a_k},\\
\text{prime-aware correction:} \quad & U^{(p_k)}_{a_k,b_k} = X^{b_k} Z^{a_k} P(\phi_{p_k}),
\end{align*}
where $P(\phi) = \mathrm{diag}(1,e^{-i\phi})$ is a $Z$-axis phase gate.

\paragraph{Lemma (Prime-Indexed Direct-Sum Teleportation).}

\begin{enumerate}
\item With prime-aware corrections $U^{(p_k)}_{a_k,b_k}$, the final logical tableau on $C_1,\dots,C_n$ is identical (up to Pauli relabeling) to the input tableau on $A_1,\dots,A_n$. All phase entries in the logical operators vanish, and the implemented channel is the identity:
\[
\mathcal{E}^{(\vec p)} = \mathrm{id}_{(\mathbb{C}^2)^{\otimes n}}.
\]
The Teleportation Axiom is thus realized in a prime-separable fashion: each prime-labeled block $p_k$ has its deformation $\phi_{p_k}$ locally cancelled.

\item With Pauli-only corrections $V_{a_k,b_k}$ and a product input $|\psi_{\mathrm{in}}\rangle = \bigotimes_{k=1}^n (\alpha_k|0\rangle+\beta_k|1\rangle)$, the output state on $C_1,\dots,C_n$ is
\[
|\psi_{\mathrm{out}}^{(\vec p)}\rangle
= \bigotimes_{k=1}^n \bigl(\alpha_k |0\rangle + \beta_k e^{i\phi_{p_k}} |1\rangle\bigr).
\]
The average teleportation fidelity over Haar-random product inputs factorizes:
\[
\overline{F}_{\mathrm{tot}}(\vec\phi)
= \prod_{k=1}^n \frac{2+\cos\phi_{p_k}}{3}.
\]
Perfect fidelity occurs if and only if each $\phi_{p_k}\equiv 0 \pmod{2\pi}$; any non-zero phase at any prime reduces the product.
\end{enumerate}

In cohomological terms, the deformation $1$-cochain is the direct sum
\[
\Phi = \bigoplus_{k=1}^n \Phi(p_k),
\]
with each $\Phi(p_k)$ supported on the edge $(B_k,C_k)$. The prime-aware correction supplies a $0$-cochain $\Psi = \bigoplus_{k=1}^n \Psi(p_k)$ such that $\delta\Psi = \Phi$ block-wise. Pauli-only correction corresponds to $\Psi = 0$, leaving a nontrivial cohomology class detected by the fidelity product above.

\subsection{Arithmetic Coherence Axiom for Prime-Indexed Teleportation}

The deformation phases $\phi_p$ can be promoted from arbitrary real parameters to values of arithmetic functors on primes. This yields a bridge between quantum teleportation and number-theoretic multiplicities.

\begin{axiom}[Arithmetic Coherence for Prime-Indexed Teleportation]
Let $\mathcal{P}$ be a set of prime labels and let
\[
A,B : \mathcal{P} \longrightarrow U(1)
\]
be arithmetic functors, assigning unit phases to primes (for example from Pell regulators or Dirichlet characters). For each $p\in\mathcal{P}$ define the resource and correction by
\[
|\Phi\rangle_A^{(p)} = \frac{1}{\sqrt{2}}\bigl(|00\rangle + A(p)\,|11\rangle\bigr)_{BC},
\qquad
P_B(p) = \begin{pmatrix} 1 & 0 \\ 0 & \overline{B(p)} \end{pmatrix}.
\]
For a finite test set $\vec p=(p_1,\dots,p_n)\subset\mathcal{P}$, construct the $n$-qubit teleportation protocol as above, using $\bigotimes_k |\Phi\rangle_A^{(p_k)}$ as resources and block-wise corrections $X^{b_k}Z^{a_k}P_B(p_k)$.

\begin{enumerate}
\item \textbf{Identity-channel condition.} The resulting channel $\mathcal{E}^{(A,B)}$ on $(\mathbb{C}^2)^{\otimes n}$ is the identity if and only if
\[
A(p) = B(p) \quad\text{for all primes } p \text{ in the tested set } \vec p.
\]
\item \textbf{Fidelity as multiplicative pseudo-distance.} If $A$ and $B$ differ on one or more tested primes, the average teleportation fidelity (Haar measure on product inputs) factorizes as
\[
\overline{F}_{\mathrm{tot}}(A,B)
= \prod_{p\in\vec p} \frac{2+\cos\Delta_p}{3},
\qquad
\Delta_p = \arg\bigl(A(p)\,\overline{B(p)}\bigr) \in (-\pi,\pi].
\]
This functional satisfies:
\begin{itemize}
\item $\overline{F}_{\mathrm{tot}}(A,B) \in (0,1]$ for any finite $\vec p$,
\item $\overline{F}_{\mathrm{tot}}(A,B) = 1$ if and only if $A(p)=B(p)$ for all $p\in\vec p$,
\item any single mismatch ($\Delta_p\neq 0$ for some $p$) strictly lowers the product.
\end{itemize}
\end{enumerate}
\end{axiom}

Thus $\overline{F}_{\mathrm{tot}}(A,B)$ acts as a prime-indexed pseudo-distance on the space of arithmetic functors, with teleportation realizing a physical test of the equality $A=B$ on the chosen prime set.

\paragraph{Canonical instantiations.}

\begin{itemize}
\item \emph{Pell–regulator functor.} For suitable primes $p$, let $\mathbb{Q}(\sqrt{p})$ be a real quadratic field with regulator $R_p = \log\varepsilon_p$, where $\varepsilon_p>1$ is a fundamental unit of norm $1$. Set
\[
A_{\mathrm{Pell}}(p) = e^{iR_p}, \qquad \phi_p = R_p \bmod 2\pi.
\]
Teleportation with resources labeled by $A_{\mathrm{Pell}}$ is perfect on a test set $\vec p$ if and only if Bob's corrections use the same phases $B=A_{\mathrm{Pell}}$ on that set; any mismatch (for example using a different discriminant or $B(p)\equiv 1$) reduces $\overline{F}_{\mathrm{tot}}(A,B)$ according to the product above.

\item \emph{Dirichlet-character functor.} Let $\chi$ be a Dirichlet character modulo $q$ and define
\[
A_\chi(p) = \chi(p) \in U(1)
\]
for primes $p\nmid q$. Teleportation is perfect on a test set $\vec p$ if and only if $B=\chi$ on that set. Using a different character $\psi$ yields
\[
\overline{F}_{\mathrm{tot}}(\chi,\psi)
= \prod_{p\in\vec p} \frac{2+\cos\bigl(\arg(\chi(p)\,\overline{\psi(p)})\bigr)}{3},
\]
a multiplicative character-theoretic discrepancy functional. For primitive characters, this product is strictly less than $1$ unless $\chi=\psi$ on the tested primes.
\end{itemize}

In cohomological terms, $A$ supplies a $1$-cochain on entangled edges, $B$ a $0$-cochain on Bob's vertex. The identity-channel condition is once again $\delta B = A$ on the tested set; the fidelity functional measures the magnitude of the residual cohomology class when this fails.

\subsection{Tele--Genesis Coupling: Coherence as a Persistence Functional}

The arithmetic coherence functional can be promoted to a dynamical quantity and coupled to a macroscopic persistence variable within a Genesis-style ODE framework.

Fix a finite test set $\vec p=(p_1,\dots,p_n)$, a reference functor $A$, and a time-dependent correction functor $B(p,t)$. Define the phase discrepancy
\[
\Delta_p(t) = \arg\bigl(A(p)\,\overline{B(p,t)}\bigr) \in (-\pi,\pi]
\]
and the prime-indexed coherence potential
\[
V_{\vec p}(t)
= \prod_{p\in\vec p} v\bigl(\Delta_p(t)\bigr),
\qquad
v(\Delta) = \frac{2+\cos\Delta}{3} \in \bigl[\tfrac13,1\bigr].
\]
The negative logarithm
\[
\Phi_{\vec p}(t) = -\log V_{\vec p}(t)
= -\sum_{p\in\vec p} \log\!\left(\frac{2+\cos\Delta_p(t)}{3}\right) \ge 0
\]
is a global mismatch potential, with $\Phi_{\vec p}(t)=0$ if and only if $\Delta_p(t)=0$ for all tested primes.

\paragraph{Arithmetic relaxation dynamics.}

A minimal gradient-type relaxation of the discrepancy field is given by
\[
\frac{d\Delta_p}{dt}
= -\eta_p \frac{\partial \Phi_{\vec p}}{\partial \Delta_p}
= \eta_p\,\frac{-\sin\Delta_p}{2+\cos\Delta_p},
\]
with mobility coefficients $\eta_p>0$. For $|\Delta_p|\ll 1$ this linearizes to
\[
\frac{d\Delta_p}{dt} \approx -\frac{\eta_p}{3}\,\Delta_p,
\]
so small discrepancies decay exponentially. The flow monotonically increases $V_{\vec p}(t)$ and drives the configuration toward the coherent diagonal $\Delta_p=0$ mod $2\pi$.

\paragraph{Coupling to a Genesis persistence variable.}

Let $C_X(t)$ be a substrate-indexed Genesis persistence variable (for example a metallurgical coherence index). If $f_X(C_X,t)$ denotes its intrinsic Genesis dynamics, Tele--Genesis coupling augments the ODE by a coherence-modulated drag term:
\[
\frac{dC_X}{dt}
= f_X(C_X,t) \;-\; \gamma_X\,\bigl(1 - V_{\vec p}(t)\bigr),
\qquad \gamma_X \ge 0.
\]
When $V_{\vec p}(t)\approx 1$ (arithmetic functors coherent on the tested primes), the additional term is negligible and $C_X(t)$ follows its intrinsic trajectory. As coherence degrades and $V_{\vec p}(t)$ decreases, the term $-\gamma_X(1-V_{\vec p})$ acts as an additional decay or impedance, an \emph{arithmetic decoherence drag} on the macroscopic persistence variable.

In this way, prime-indexed teleportation serves as a dynamical probe of arithmetic coherence within the Genesis architecture: the same prime-labeled multiplicity structure underlies both the stabilizer tableau decomposition and the persistence-modulated ODE, with the teleportation fidelity $V_{\vec p}(t)$ functioning as a shared scalar observable linking quantum information, arithmetic functors, and nonlinear persistence dynamics.

\appendix
\section*{Mathematical Appendix}

\section{Operator Definitions and Domains}

\subsection{Prime Sector Operator \texorpdfstring{$\mathbf{A}_{\rm prime}$}{A\_prime}}

Let \(\mathbb{P}\) denote the set of prime numbers and consider the Hilbert space
\[
  \ell^2(\mathbb{P})
  = \Bigl\{ x = (x_p)_{p\in\mathbb{P}} : \sum_{p\in\mathbb{P}} |x_p|^2 < \infty \Bigr\},
\]
with canonical orthonormal basis \(\{e_p : p\in\mathbb{P}\}\), where \(e_p\) has \(1\) in the \(p\)-th coordinate and \(0\) elsewhere.

Let \(\sigma:\mathbb{P}\to\mathbb{R}\) be a bounded function. Define the multiplication operator
\[
  D_\sigma : \ell^2(\mathbb{P}) \to \ell^2(\mathbb{P}), \qquad D_\sigma e_p = \sigma(p) e_p.
\]
Then \(D_\sigma\) is bounded and self–adjoint with
\[
  \|D_\sigma\| = \sup_{p\in\mathbb{P}} |\sigma(p)|.
\]

Let \(K:\ell^2(\mathbb{P})\to\ell^2(\mathbb{P})\) be a Hilbert–Schmidt operator, i.e.
\[
  \|K\|_{\rm HS}^2 = \sum_{p,q\in\mathbb{P}} |k_{pq}|^2 < \infty,
\]
where \(k_{pq} = \langle e_p, K e_q \rangle\). Then \(K\) is compact and bounded with
\[
  \|K\| \le \|K\|_{\rm HS}.
\]

Define
\[
  \mathbf{A}_{\rm prime} := D_\sigma + K.
\]
Since \(\sigma\) is bounded and real–valued, \(D_\sigma\) is self–adjoint, and if \(K\) is self–adjoint then \(\mathbf{A}_{\rm prime}\) is self–adjoint on the same domain, with
\[
  \|\mathbf{A}_{\rm prime}\| \le \|D_\sigma\| + \|K\|
  \le \sup_{p}|\sigma(p)| + \|K\|_{\rm HS}.
\]

Moreover, by Weyl’s theorem on the invariance of the essential spectrum under compact perturbations,
\[
  \sigma_{\rm ess}(\mathbf{A}_{\rm prime})
  = \sigma_{\rm ess}(D_\sigma)
  = \overline{\sigma(\mathbb{P})}^{\,\rm ess},
\]
the essential closure of the image of \(\sigma\), since \(K\) is compact.

\subsection{Zeta Sector Operator \texorpdfstring{$\mathbf{B}_{\rm zeta}$}{B\_zeta}}

Let \(H_\zeta = L^2(\mathbb{R})\) with Fourier transform
\[
  \hat{f}(\xi) = \frac{1}{\sqrt{2\pi}}\int_{\mathbb{R}} e^{-ix\xi} f(x)\,dx.
\]
Let \(m:\mathbb{R} \to \mathbb{R}\) be a bounded measurable function derived from the Riemann zeta function; e.g.\ for some fixed \(\sigma_0\) and mollifier \(w\),
\[
  m(\xi) = \int_{\mathbb{R}} w(\xi-\eta)\, \Re\log\zeta\bigl(\tfrac12 + i\eta\bigr)\,d\eta.
\]
Assume \(m\in L^\infty(\mathbb{R})\) and real–valued almost everywhere. Define
\[
  (\widehat{\mathbf{B}_{\rm zeta} f})(\xi) := m(\xi) \hat{f}(\xi).
\]
Then \(\mathbf{B}_{\rm zeta}\) is bounded and self–adjoint on \(H_\zeta\), with
\[
  \|\mathbf{B}_{\rm zeta}\| = \|m\|_{L^\infty}.
\]
The essential spectrum satisfies
\[
  \sigma_{\rm ess}(\mathbf{B}_{\rm zeta})
  = \operatorname{EssRan}(m),
\]
the essential range of \(m\).

\subsection{Internal Sector \texorpdfstring{$\mathbf{E}_{\rm internal}$}{E\_internal}}

Let \(H_{\rm int} = \mathbb{C}^d\) and let \(\mathbf{E}_{\rm internal}\) be any Hermitian matrix in \(\mathbb{C}^{d\times d}\). Then \(\mathbf{E}_{\rm internal}\) is bounded and self–adjoint with finite spectrum.

\subsection{Bridge Operator \texorpdfstring{$\mathbf{\Xi}_{\rm bridge}$}{Xi\_bridge}}

Consider the composite Hilbert space
\[
  H = \ell^2(\mathbb{P}) \,\widehat{\otimes}\, L^2(\mathbb{R}) \,\widehat{\otimes}\, \mathbb{C}^d.
\]
For each prime \(p_i\) and each nontrivial zero \(\rho_k = \tfrac12 + i\gamma_k\) of \(\zeta(s)\), define a rank–one operator \(\Pi_{i,k}\) on \(H\) projecting onto the mode labeled by \((p_i,\rho_k)\) and some internal basis vector.

Let \(\kappa:\mathbb{P}\times\{\rho_k\} \to \mathbb{R}\) be a coupling kernel, and let \(\varphi\) be a compactly supported test function on \(\mathbb{R}\) (e.g.\ Fejér kernel in \(\gamma\)). Define
\[
  \mathbf{\Xi}_{\rm bridge}
  := \sum_{i,k} \kappa(p_i,\rho_k) \, \varphi(\gamma_k/T)\, \cos(\gamma_k\log p_i) \,\Pi_{i,k},
\]
with the sum converging in operator norm by requiring
\[
  \sum_{i,k} |\kappa(p_i,\rho_k)|^2 |\varphi(\gamma_k/T)|^2 < \infty.
\]
Then \(\mathbf{\Xi}_{\rm bridge}\) is Hilbert–Schmidt and hence compact; moreover, if all coefficients are real and \(\Pi_{i,k}\) are self–adjoint, then \(\mathbf{\Xi}_{\rm bridge}\) is self–adjoint.

\section{Zeta–Multiplicity Operator and Norm Bounds}

\subsection{Definition and Self–Adjointness}

Define the full operator
\[
  \mathcal{U}_{\rm ZMO}
  := \mathbf{A}_{\rm prime} \otimes I_{H_\zeta\otimes H_{\rm int}}
     + I_{\ell^2(\mathbb{P})}\otimes \mathbf{B}_{\rm zeta} \otimes I_{H_{\rm int}}
     + I_{\ell^2(\mathbb{P})\otimes H_\zeta}\otimes \mathbf{E}_{\rm internal}
     + \Lambda_m\, \mathbf{\Xi}_{\rm bridge},
\]
acting on \(H = \ell^2(\mathbb{P})\widehat{\otimes}L^2(\mathbb{R})\widehat{\otimes}\mathbb{C}^d\), with \(\Lambda_m\in(0,1)\) a scalar.

Each tensor–extended operator is bounded and self–adjoint. If \(\mathbf{\Xi}_{\rm bridge}\) is self–adjoint, then \(\Lambda_m \mathbf{\Xi}_{\rm bridge}\) is self–adjoint. Thus \(\mathcal{U}_{\rm ZMO}\) is bounded and self–adjoint on all of \(H\), with operator norm bounded by
\begin{align*}
  \|\mathcal{U}_{\rm ZMO}\|
  &\le \|\mathbf{A}_{\rm prime}\| + \|\mathbf{B}_{\rm zeta}\| + \|\mathbf{E}_{\rm internal}\|
       + |\Lambda_m|\,\|\mathbf{\Xi}_{\rm bridge}\| \\
  &\le \sup_p|\sigma(p)| + \|K\|_{\rm HS} + \|m\|_{L^\infty}
       + \|\mathbf{E}_{\rm internal}\|
       + |\Lambda_m|\,\|\mathbf{\Xi}_{\rm bridge}\|_{\rm HS}.
\end{align*}

\subsection{Essential Spectrum}

Since \(\mathbf{\Xi}_{\rm bridge}\) is compact, so is \(\Lambda_m\mathbf{\Xi}_{\rm bridge}\). Thus \(\mathcal{U}_{\rm ZMO}\) is a compact perturbation of
\[
  \mathcal{U}_0
  := \mathbf{A}_{\rm prime} \otimes I_{H_\zeta\otimes H_{\rm int}}
     + I_{\ell^2(\mathbb{P})}\otimes \mathbf{B}_{\rm zeta} \otimes I_{H_{\rm int}}
     + I_{\ell^2(\mathbb{P})\otimes H_\zeta}\otimes \mathbf{E}_{\rm internal}.
\]
By the stability of essential spectrum under compact perturbations,
\[
  \sigma_{\rm ess}(\mathcal{U}_{\rm ZMO})
  = \sigma_{\rm ess}(\mathcal{U}_0).
\]
Since \(\mathcal{U}_0\) is a sum of commuting self–adjoint operators acting on tensor factors, standard results on spectra of sums of commuting normal operators yield
\[
  \sigma_{\rm ess}(\mathcal{U}_0)
  = \sigma_{\rm ess}(\mathbf{A}_{\rm prime})
    + \sigma_{\rm ess}(\mathbf{B}_{\rm zeta})
    + \sigma_{\rm ess}(\mathbf{E}_{\rm internal}),
\]
where the sum on the right is the Minkowski sum of the sets.\medskip

Given
\[
  \sigma_{\rm ess}(\mathbf{A}_{\rm prime})
  = \overline{\sigma(\mathbb{P})}^{\,\rm ess}, \qquad
  \sigma_{\rm ess}(\mathbf{B}_{\rm zeta}) = \operatorname{EssRan}(m), \qquad
  \sigma_{\rm ess}(\mathbf{E}_{\rm internal}) = \sigma(\mathbf{E}_{\rm internal}),
\]
we obtain the explicit essential spectrum description:
\[
  \sigma_{\rm ess}(\mathcal{U}_{\rm ZMO})
  = \overline{\sigma(\mathbb{P})}^{\,\rm ess}
    + \operatorname{EssRan}(m)
    + \sigma(\mathbf{E}_{\rm internal}).
\]

\subsection{Unitary and Contractive Evolutions}

Since \(\mathcal{U}_{\rm ZMO}\) is bounded and self–adjoint, it generates a strongly continuous unitary group
\[
  U(t) = e^{it\mathcal{U}_{\rm ZMO}}, \qquad t\in\mathbb{R}.
\]
If instead we define a generator
\[
  \mathcal{A} = -\mathcal{U}_{\rm ZMO},
\]
and view the evolution in the semigroup picture
\[
  T(t) = e^{t\mathcal{A}}, \qquad t\ge 0,
\]
then
\[
  \|T(t)\|
  = \|e^{-t\mathcal{U}_{\rm ZMO}}\|
  = e^{-t\lambda_{\min}},
\]
where \(\lambda_{\min}\) is the minimal point of \(\sigma(\mathcal{U}_{\rm ZMO})\). Imposing a spectral lower bound
\[
  \sigma(\mathcal{U}_{\rm ZMO}) \subset [\lambda_0,\infty), \qquad \lambda_0 > 0,
\]
yields a contractive semigroup with
\[
  \|T(t)\| \le e^{-t\lambda_0} \le 1 - \varepsilon_t
\]
for some \(\varepsilon_t>0\), providing explicit contractivity estimates.

A Cayley transform
\[
  C = (\mathcal{U}_{\rm ZMO}-iI)(\mathcal{U}_{\rm ZMO}+iI)^{-1}
\]
yields a unitary operator whose spectrum lies on the unit circle; the norm bound
\(\|C\|=1\) is immediate, and compact perturbations of \(\mathcal{U}_{\rm ZMO}\) translate into compact perturbations of \(C\) in the normal category.

\section{Two–Qubit Hamiltonian: Hermiticity and Bounds}

\subsection{Local Hamiltonian Hermiticity}

For each time \(t\), the local unitaries
\[
  U_{\rm dop}(t) = R(\theta_{\rm dop}(t))\otimes I_2, \qquad
  U_{\rm ser}(t) = I_2 \otimes R(\theta_{\rm ser}(t)),
\]
with
\[
  R(\theta) = \begin{pmatrix}
    \cos\theta & -i\sin\theta \\
    -i\sin\theta & \cos\theta
  \end{pmatrix}
\]
satisfy \(R(\theta)^\dagger R(\theta) = I_2\).

Define
\[
  H_{\rm local}(t)
  = -\frac{i}{2}\bigl(U_{\rm dop}(t) - U_{\rm dop}(t)^\dagger\bigr)
    -\frac{i}{2}\bigl(U_{\rm ser}(t) - U_{\rm ser}(t)^\dagger\bigr).
\]
Then
\[
  H_{\rm local}(t)^\dagger
  = \frac{i}{2}\bigl(U_{\rm dop}(t)^\dagger - U_{\rm dop}(t)\bigr)
    + \frac{i}{2}\bigl(U_{\rm ser}(t)^\dagger - U_{\rm ser}(t)\bigr)
  = H_{\rm local}(t),
\]
so \(H_{\rm local}(t)\) is Hermitian for all \(t\).

Moreover, since \(\|U_{\rm dop}(t)\| = \|U_{\rm ser}(t)\| = 1\), we have
\[
  \|H_{\rm local}(t)\|
  \le \frac12\bigl(\|U_{\rm dop}(t)\| + \|U_{\rm dop}(t)^\dagger\|
                 + \|U_{\rm ser}(t)\| + \|U_{\rm ser}(t)^\dagger\|\bigr)
  = 2.
\]

\subsection{Interaction Hamiltonian Hermiticity and Bound}

The ZZ interaction is
\[
  ZZ = Z\otimes Z = \operatorname{diag}(1,-1,-1,1),
\]
which is Hermitian with \(\|ZZ\| = 1\). For each time \(t\),
\[
  H_{\rm int}^{\rm prime}(t)
  = \Lambda_m^{\mathrm{op}}(t)\,ZZ
\]
is Hermitian because \(\Lambda_m^{\mathrm{op}}(t)\in\mathbb{R}\), and satisfies
\[
  \|H_{\rm int}^{\rm prime}(t)\|
  = |\Lambda_m^{\mathrm{op}}(t)|\,\|ZZ\|
  \le \frac12,
\]
by the construction of \(\Lambda_m^{\mathrm{op}}(t)\).

Similarly, the scalar control interaction
\[
  H_{\rm int}^{\rm scalar}(t)
  = \overline{\Lambda_m^{\mathrm{op}}}\, ZZ
\]
is Hermitian with
\[
  \|H_{\rm int}^{\rm scalar}(t)\|
  = |\overline{\Lambda_m^{\mathrm{op}}}|\,\|ZZ\|
  \le \frac12.
\]

\subsection{Total Hamiltonian and Exponential Bounds}

The total Hamiltonians
\[
  H^{\rm prime}(t) = H_{\rm local}(t) + H_{\rm int}^{\rm prime}(t), \qquad
  H^{\rm scalar}(t) = H_{\rm local}(t) + H_{\rm int}^{\rm scalar}(t)
\]
are Hermitian as sums of Hermitian operators.

The norms satisfy
\[
  \|H^{(\cdot)}(t)\|
  \le \|H_{\rm local}(t)\| + \|H_{\rm int}^{(\cdot)}(t)\|
  \le 2 + \tfrac12 = \frac{5}{2},
\]
uniformly in \(t\).

For a time step \(\Delta t\), the evolution operators
\[
  U^{(\cdot)}(t,\Delta t)
  = e^{-i H^{(\cdot)}(t)\Delta t}
\]
are unitary with \(\|U^{(\cdot)}(t,\Delta t)\|=1\). The series expansion
\[
  e^{-iH^{(\cdot)}(t)\Delta t}
  = I - i H^{(\cdot)}(t)\Delta t + \mathcal{O}(\Delta t^2)
\]
admits the operator norm bound
\[
  \|I - e^{-iH^{(\cdot)}(t)\Delta t}\|
  \le \|H^{(\cdot)}(t)\|\,\Delta t\, e^{\|H^{(\cdot)}(t)\|\Delta t}
  \le \frac{5}{2}\Delta t\, e^{\frac{5}{2}\Delta t},
\]
which controls the deviation from identity for small \(\Delta t\) and ensures numerical stability in the kernel.

\section{Concurrence and Invariance Bounds}

Given a normalized pure state \(|\psi\rangle = a|00\rangle + b|01\rangle + c|10\rangle + d|11\rangle \in \mathbb{C}^2\otimes\mathbb{C}^2\), the concurrence satisfies
\[
  0 \le C(\psi) = 2|ad - bc| \le 1.
\]
The upper bound follows from the inequality
\[
  |ad - bc|
  \le \sqrt{(|a|^2 + |b|^2)(|c|^2 + |d|^2)}
  \le \frac12\bigl(|a|^2+|b|^2+|c|^2+|d|^2\bigr)
  = \frac12,
\]
so \(C(\psi) \le 1\).

Under any global unitary \(U\in U(4)\), the concurrence is invariant:
\[
  C(U\psi) = C(\psi),
\]
since concurrence depends only on the Schmidt coefficients of \(|\psi\rangle\), which are unchanged by local or global unitaries. In particular, the Bell state \((|00\rangle+|11\rangle)/\sqrt{2}\) has \(C=1\).

In the discrete–time kernel, the numerical scheme is built from exact exponentials of Hermitian matrices, so each step is a unitary transform up to floating–point error. Thus the theoretical concurrence is preserved; the observed deviations in the matched–control experiment arise from the difference in the continuous–time dynamics under \(H^{\rm prime}(t)\) and \(H^{\rm scalar}(t)\), not from non–unitarity.

\appendix
\section{Mathematical Appendix: Prime–Indexed Teleportation Proofs and Bounds}

This appendix collects explicit derivations and norm estimates underlying the prime–indexed teleportation constructions used in the main text. Throughout, $\mathcal{H}\cong\mathbb{C}^2$ denotes the single–qubit Hilbert space, and $\|\cdot\|_\infty$ is the operator norm on $\mathcal{B}((\mathbb{C}^2)^{\otimes n})$.

\subsection{Single–Qubit Prime–Indexed Teleportation}

Fix a prime label $p$ and a phase $\phi_p\in\mathbb{R}/2\pi\mathbb{Z}$. Consider the deformed Bell pair
\begin{equation}
  \label{eq:phi-p-state}
  \ket{\Phi}^{(p)}_{BC}
  = \frac{1}{\sqrt{2}}\bigl(\ket{00} + e^{i\phi_p}\ket{11}\bigr)_{BC}.
\end{equation}
Let Alice hold an input qubit $A$ in an arbitrary pure state
\[
  \ket{\psi}_{A} = \alpha\ket{0} + \beta\ket{1}, \qquad |\alpha|^2+|\beta|^2=1,
\]
and the resource pair $BC$ be in the state \eqref{eq:phi-p-state}. The total initial state is
\[
  \ket{\Psi_{\text{in}}}
  = \ket{\psi}_A \otimes \ket{\Phi}^{(p)}_{BC}
  = \frac{1}{\sqrt{2}}
    \bigl(\alpha\ket{0} + \beta\ket{1}\bigr)_A
    \otimes
    \bigl(\ket{00} + e^{i\phi_p}\ket{11}\bigr)_{BC}.
\]

\subsubsection{State Evolution Under the Teleportation Circuit}

Alice applies CNOT$_{A\to B}$ followed by $H_A$, then measures $(A,B)$ in the computational basis.

\paragraph{CNOT step.}
Applying CNOT$_{A\to B}$ gives
\begin{align*}
  \ket{\Psi_1}
  &= \frac{1}{\sqrt{2}}
     \Bigl[
       \alpha\ket{0}_A\ket{00}_{BC}
       + \alpha e^{i\phi_p}\ket{0}_A\ket{11}_{BC}
       + \beta\ket{1}_A\ket{10}_{BC}
       + \beta e^{i\phi_p}\ket{1}_A\ket{01}_{BC}
     \Bigr].
\end{align*}

\paragraph{Hadamard on $A$.}
Using $H\ket{0} = (\ket{0}+\ket{1})/\sqrt{2}$ and $H\ket{1} = (\ket{0}-\ket{1})/\sqrt{2}$,
\begin{align*}
  \ket{\Psi_2}
  &= (H_A\otimes I_{BC})\ket{\Psi_1} \\
  &= \frac{1}{2}
     \Bigl[
       \alpha (\ket{0}+\ket{1})_A \ket{00}_{BC}
       + \alpha e^{i\phi_p} (\ket{0}+\ket{1})_A \ket{11}_{BC} \\
  &\qquad\quad
       + \beta (\ket{0}-\ket{1})_A \ket{10}_{BC}
       + \beta e^{i\phi_p} (\ket{0}-\ket{1})_A \ket{01}_{BC}
     \Bigr].
\end{align*}
This can be regrouped by measurement outcomes $(a,b)\in\{0,1\}^2$ on $(A,B)$:
\begin{align*}
  \ket{\Psi_2}
  &= \frac{1}{2}\Bigl(
     \ket{00}_{AB}\bigl[\alpha\ket{0} + \beta e^{i\phi_p}\ket{1}\bigr]_C
   + \ket{01}_{AB}\bigl[\alpha e^{i\phi_p}\ket{1} + \beta\ket{0}\bigr]_C \\
  &\hspace{3.3em}
   + \ket{10}_{AB}\bigl[\alpha\ket{0} - \beta e^{i\phi_p}\ket{1}\bigr]_C
   + \ket{11}_{AB}\bigl[\alpha e^{i\phi_p}\ket{1} - \beta\ket{0}\bigr]_C
   \Bigr).
\end{align*}

Thus, conditional on Alice obtaining outcome $(a,b)$, Bob's unnormalized post–measurement state on $C$ is
\[
  \ket{\psi_{ab}}_C \propto X^b Z^a \bigl(\alpha\ket{0} + \beta e^{i\phi_p}\ket{1}\bigr)_C.
\]

\subsubsection{Pauli–Only Correction and Residual Error}

If Bob applies the usual Pauli correction $X^b Z^a$, the remaining state is
\[
  \ket{\psi_{\text{out}}^{(p)}}
  = \alpha\ket{0} + \beta e^{i\phi_p}\ket{1}.
\]
So the implemented channel is
\[
  \mathcal{E}^{(p)}_{\text{Pauli}}(\rho)
  = U_{\phi_p} \, \rho \, U_{\phi_p}^\dagger,
  \qquad
  U_{\phi_p} = \ket{0}\!\bra{0} + e^{i\phi_p}\ket{1}\!\bra{1},
\]
a $Z$–axis phase rotation.

\subsubsection{Prime–Aware Correction and Exact Identity}

If Bob instead appends a prime–aware $Z$–phase gate
\[
  P(\phi_p) = \mathrm{diag}(1,e^{-i\phi_p}),
\]
then
\[
  P(\phi_p)\ket{\psi_{\text{out}}^{(p)}}
  = \alpha\ket{0} + \beta e^{i\phi_p} e^{-i\phi_p}\ket{1}
  = \alpha\ket{0} + \beta\ket{1}
  = \ket{\psi}.
\]
Hence the combined channel with prime–aware corrections
\[
  \mathcal{E}^{(p)}_{\text{prime-aware}} = \mathrm{id}_{\mathcal{B}(\mathcal{H})}.
\]

\subsubsection{Single–Qubit Fidelity and Norm Bounds}

Let $\Lambda_p$ denote the Pauli–only teleportation channel with resource phase $\phi_p$. It is a unitary channel induced by $U_{\phi_p}$. For a pure input $\ket{\psi}$, the fidelity is
\[
  F_p(\ket{\psi})
  = \bigl|\langle\psi|U_{\phi_p}|\psi\rangle\bigr|^2
  = \bigl|\;|\alpha|^2 + e^{i\phi_p}|\beta|^2\;\bigr|^2
  = |\alpha|^4 + |\beta|^4 + 2|\alpha|^2|\beta|^2 \cos\phi_p.
\]
Averaging over the Haar measure on single–qubit pure states yields
\begin{equation}
  \label{eq:single-qubit-avg-fidelity}
  \overline{F}_p
  = \frac{2+\cos\phi_p}{3}.
\end{equation}
This is the well–known single–qubit relation between average fidelity and Bloch–sphere contraction for a unitary phase error.

The diamond–norm distance between $\Lambda_p$ and the identity channel satisfies
\[
  \|\Lambda_p - \mathrm{id}\|_\diamond
  = 2\sup_{\ket{\psi}} \sqrt{1 - F_p(\ket{\psi})}
  \le 2\sqrt{1 - \min_{\ket{\psi}} F_p(\ket{\psi})}.
\]
The minimum fidelity over pure $\ket{\psi}$ occurs for $|\alpha|^2=\tfrac12$, yielding
\[
  \min_{\ket{\psi}} F_p(\ket{\psi})
  = \frac{1+\cos\phi_p}{2},
\]
so
\begin{equation}
  \label{eq:single-qubit-diamond-bound}
  \|\Lambda_p - \mathrm{id}\|_\diamond
  \le 2\sqrt{1 - \frac{1+\cos\phi_p}{2}}
  = 2\sqrt{\frac{1-\cos\phi_p}{2}}
  = 2\left|\sin\frac{\phi_p}{2}\right|.
\end{equation}
In particular, as $\phi_p\to 0$ this behaves as $\|\Lambda_p - \mathrm{id}\|_\diamond \le |\phi_p| + O(\phi_p^3)$.

\subsection{\texorpdfstring{$n$}{n}–Qubit Direct–Sum Proofs}

Let $\vec p=(p_1,\dots,p_n)$ and consider the resource
\[
  \ket{\Phi}^{(\vec p)}
  = \bigotimes_{k=1}^n \ket{\Phi}^{(p_k)}_{B_kC_k},
\]
with each block as in \eqref{eq:phi-p-state}. Let Alice's input be the product state
\[
  \ket{\psi_{\mathrm{in}}}
  = \bigotimes_{k=1}^n (\alpha_k\ket{0}+\beta_k\ket{1})_{A_k}.
\]
The protocol applies the single–qubit teleportation to each triple $(A_k,B_k,C_k)$ independently.

\subsubsection{Exact Factorization of the Channel}

Because the gate pattern is a tensor product of local CNOT–Hadamard–measurement–correction routines, the overall channel factors as the tensor product of the single–qubit channels:
\[
  \Lambda_{\vec p}
  = \bigotimes_{k=1}^n \Lambda_{p_k},
\]
with $\Lambda_{p_k}$ the Pauli–only one–qubit channel induced by $U_{\phi_{p_k}}$.

\paragraph{Prime–Aware case.}
With prime–aware corrections, each $\Lambda_{p_k}$ is the identity, hence
\[
  \Lambda_{\vec p}^{\text{prime-aware}}
  = \bigotimes_{k=1}^n \mathrm{id}
  = \mathrm{id}_{(\mathbb{C}^2)^{\otimes n}}.
\]

\paragraph{Pauli–only case.}
With Pauli corrections only, each block yields
\[
  \ket{\psi^{(k)}_{\text{out}}}
  = \alpha_k\ket{0} + \beta_k e^{i\phi_{p_k}}\ket{1},
\]
so the total output is
\[
  \ket{\psi_{\text{out}}^{(\vec p)}}
  = \bigotimes_{k=1}^n
    \bigl(\alpha_k\ket{0} + \beta_k e^{i\phi_{p_k}}\ket{1}\bigr).
\]

\subsubsection{Average Fidelity Factorization}

The pure–state fidelity between $\ket{\psi_{\mathrm{in}}}$ and $\ket{\psi_{\mathrm{out}}^{(\vec p)}}$ is
\begin{align*}
  F_{\vec p}(\{\alpha_k,\beta_k\})
  &= \bigl|\langle\psi_{\mathrm{in}}|\psi_{\mathrm{out}}^{(\vec p)}\rangle\bigr|^2 \\
  &= \prod_{k=1}^n
     \bigl|\;|\alpha_k|^2 + e^{i\phi_{p_k}}|\beta_k|^2\;\bigr|^2 \\
  &= \prod_{k=1}^n
     \Bigl(|\alpha_k|^4 + |\beta_k|^4
           + 2|\alpha_k|^2|\beta_k|^2\cos\phi_{p_k}\Bigr).
\end{align*}
Averaging independently over Haar–random single–qubit pure states yields
\begin{align*}
  \overline{F}_{\mathrm{tot}}(\vec\phi)
  &= \prod_{k=1}^n \overline{F}_{p_k}
   = \prod_{k=1}^n \frac{2+\cos\phi_{p_k}}{3}.
\end{align*}
This establishes the multiplicative form used in the main text.

\subsubsection{Operator Norm Bounds in the Multi–Prime Case}

Let $\Lambda_{\vec p}$ denote the Pauli–only $n$–qubit teleportation channel and $\mathbb{I}$ the identity channel. Because $\Lambda_{\vec p}$ is a tensor product of single–qubit unitaries, its diamond–norm distance from $\mathbb{I}$ obeys
\[
  \|\Lambda_{\vec p}-\mathbb{I}\|_\diamond
  = \left\|\bigotimes_{k=1}^n \Lambda_{p_k} - \bigotimes_{k=1}^n \mathrm{id}\right\|_\diamond.
\]
Using subadditivity and the product expansion
\[
  \bigotimes_{k=1}^n \Lambda_{p_k} - \bigotimes_{k=1}^n \mathrm{id}
  = \sum_{\emptyset\neq S\subseteq\{1,\dots,n\}}
      \bigotimes_{k\in S}(\Lambda_{p_k}-\mathrm{id})
      \bigotimes_{j\notin S} \mathrm{id},
\]
and the stability of the diamond norm under tensoring with identities, we obtain
\begin{align*}
  \|\Lambda_{\vec p}-\mathbb{I}\|_\diamond
  &\le \sum_{\emptyset\neq S}
        \prod_{k\in S} \|\Lambda_{p_k}-\mathrm{id}\|_\diamond.
\end{align*}
Using the single–qubit bound \eqref{eq:single-qubit-diamond-bound},
\begin{equation}
  \label{eq:n-qubit-diamond-bound}
  \|\Lambda_{\vec p}-\mathbb{I}\|_\diamond
  \le \sum_{\emptyset\neq S\subseteq\{1,\dots,n\}}
       \prod_{k\in S} 2\left|\sin\frac{\phi_{p_k}}{2}\right|.
\end{equation}
In the small–angle regime $|\phi_{p_k}|\ll 1$, this provides the estimate
\[
  \|\Lambda_{\vec p}-\mathbb{I}\|_\diamond
  \le 2\sum_{k=1}^n \left|\sin\frac{\phi_{p_k}}{2}\right|
    + O\!\left(\sum_{k\neq j}|\phi_{p_k}\phi_{p_j}|\right)
  \approx \sum_{k=1}^n |\phi_{p_k}| + O(\|\vec\phi\|_2^2).
\]

\subsection{Arithmetic Functors and Coherence Functional}

Let $A,B:\mathcal{P}\to U(1)$ be arithmetic functors and $\vec p=(p_1,\dots,p_n)$ a finite test set of primes. Define
\[
  \Delta_p = \arg\bigl(A(p)\,\overline{B(p)}\bigr) \in (-\pi,\pi], \qquad
  v(\Delta) = \frac{2+\cos\Delta}{3}.
\]
Then the prime–indexed coherence functional is
\[
  V_{\vec p}(A,B) = \prod_{p\in\vec p} v(\Delta_p),
\]
and the associated mismatch potential
\[
  \Phi_{\vec p}(A,B)
  = -\log V_{\vec p}(A,B)
  = -\sum_{p\in\vec p} \log v(\Delta_p) \ge 0.
\]

\subsubsection{Monotonicity Under Gradient Flow}

Consider the ODE
\[
  \frac{d\Delta_p}{dt}
  = -\eta_p \frac{\partial \Phi_{\vec p}}{\partial \Delta_p}
  = \eta_p\,\frac{-\sin\Delta_p}{2+\cos\Delta_p},
  \qquad \eta_p>0.
\]
Then
\begin{align*}
  \frac{d}{dt}\Phi_{\vec p}
  &= \sum_{p\in\vec p}
     \frac{\partial\Phi_{\vec p}}{\partial\Delta_p}
     \frac{d\Delta_p}{dt} \\
  &= -\sum_{p\in\vec p}
        \frac{\partial\Phi_{\vec p}}{\partial\Delta_p}
        \eta_p \frac{\partial\Phi_{\vec p}}{\partial\Delta_p}
   = -\sum_{p\in\vec p} \eta_p
        \left(\frac{\partial\Phi_{\vec p}}{\partial\Delta_p}\right)^2
  \le 0.
\end{align*}
Thus $\Phi_{\vec p}(t)$ is nonincreasing along the flow, and $V_{\vec p}(t)=e^{-\Phi_{\vec p}(t)}$ is nondecreasing.

The critical points satisfy $\sin\Delta_p=0$ for all $p$, i.e.\ $\Delta_p\in\{0,\pi\}$. One checks that $\Delta_p=0$ is stable (local minimum of $\Phi_{\vec p}$) and $\Delta_p=\pi$ is unstable (local maximum for the $p$–component).

\subsection{Coupling to a Persistence Variable}

Let $C_X(t)$ be a real–valued persistence observable for a substrate $X$, evolving intrinsically by
\[
  \frac{dC_X}{dt} = f_X(C_X,t).
\]
Tele–Genesis coupling augments this by
\[
  \frac{dC_X}{dt}
  = f_X(C_X,t) - \gamma_X\bigl(1 - V_{\vec p}(t)\bigr),
  \qquad \gamma_X\ge 0,
\]
with $V_{\vec p}(t)=V_{\vec p}(A,B(t))$ as above.

\subsubsection{Bounds and Qualitative Behaviour}

Since $v(\Delta)\in[\tfrac13,1]$, we have
\[
  \frac{1}{3^n}
  \le V_{\vec p}(t) \le 1,
\]
so the additional term is bounded:
\[
  -\gamma_X \le -\gamma_X(1-V_{\vec p}(t)) \le 0.
\]
In particular, the Tele–Genesis coupling never injects positive drift into $C_X$; it only contributes an additional nonpositive ``drag'' bounded in magnitude by $\gamma_X$.

If $f_X$ admits a stable fixed point $C_X^\star$ in the absence of coupling, small $\gamma_X$ yields a perturbation of that fixed point and its basin; in the limit $\gamma_X\to 0$ one recovers the uncoupled Genesis dynamics. For fixed $\gamma_X>0$, as the arithmetic flow drives $V_{\vec p}(t)\to 1$, the perturbative term vanishes asymptotically and the long–time behaviour of $C_X$ returns to that of the intrinsic dynamics.

These estimates justify the interpretation of $V_{\vec p}(t)$ as a coherence–controlled persistence functional: arithmetic mismatch (low $V_{\vec p}$) temporarily enhances effective decay or impedance for $C_X$, while restoration of arithmetic coherence (high $V_{\vec p}$) removes that drag.

\nocite{*}
\bibliographystyle{plain}
\bibliography{references}

\end{document}

